\emph{Design of the AIT replay.} Per organisation, \cref{apptab:aitsupp} trains a
leave-one-organisation-out detector whose flow features are chosen per fold from the training
organisations only, calibrates on in-organisation benign flows strictly preceding the first attack,
forms $(\mathrm{SrcIP},\mathrm{DstIP},\text{2\,h})$ episodes, runs e-LOND, and suppresses each detected
malicious episode by replaying real benign flows to an attacked victim of the organisation (200
draws). A firing pad contributes the maximum e-value and raises the cost, which the zero-pad closed form
misses. The two arms replay differently because the two detectors read different things. Flow-only:
flow features carry no context, so pads are drawn with replacement from the pool's observed
$0$-or-$\ceil$ evidence, an empirical i.i.d.\ generator, so a modelled cost may exceed the finite
observed pool. Host-conditioned: the six host features are causal, so appending pads changes the
context every later pad is scored in, and holding them fixed would be a static-context diagnostic
rather than a replay; we therefore recompute the victim's context as pads accumulate (counts raised,
distinct-peer counts pinned, failure fractions moved to the pool's rate), score the whole pool under
it, and report a success only over the accumulation range actually scored. The chain runs end-to-end
at all eight organisations, in both arms: under first-flow order 84 of 85 flow-only detections are
suppressible, per-draw success $0.79$--$1.00$, and 57 of 63 host-conditioned
detections are, per-draw success $0.61$--$1.00$. Host conditioning therefore blunts the attack
without stopping it. The causal replay is what shows both halves. It is why most episodes still
suppress: \texttt{shaw}'s ordinary-to-victim flows fire 7\% in
their own contexts, which under a static-context replay made no episode suppressible, but a pad's
context is the attacked pair's, and under that context the same pool fires at most
0.5\% across accumulation, so all three episodes suppress. It is also why the
remaining 6 do not (\texttt{harrison}, \texttt{wardbeck}, \texttt{wheeler}): under the accumulated context a pad can
itself fire, and a firing pad adds the ceiling to the group's evidence instead of diluting it, so on
those episodes no draw brings the group mean below threshold in any of the 200 draws, within twice
the zero-pad cost (at least $1{,}000$ pads), which is the range the context curve is scored over. What
decides the outcome is the pad count an episode needs against the rate at which pads fire along that
accumulation path. These are censored simulations rather than proof that the episodes cannot be
padded: a budget or a draw count large enough might still find a route. The pinning
models a single already-active attacker--victim pair; padding from new hosts would move the
distinct-peer counts and is not covered.
